\chapter{Integration Notes}

\section{Integration into a Project}
\ifnum \DocumentGroup=1{
Please refer to the installation instructions included in the RTA-SUM Getting Started Guide to generate the RTA-SUM SSFM directly inside ISOLAR-AB.
Briefly you can select RTA Code Generator for configuring the SUM Modules that want to be integrated in the project accordingly to its SDV 2.0 version. 
When the option "Not Configured" is selected, the module will not be integrated.}
\else
{
There is no need for standard installation procedure for SSFM. The folder itself contains all the necessary
artifacts to integrate into a project. Dropping the complete folder structure into a BSW configuration
project will enable the user to configure SSFM and generate all the necessary artifacts.
}
\fi

Refer to Chapter~\ref {ch:IntGuide}
for more details about how to integrate this module into a project.

\section{Module Dependencies}
\label{sec:dependencies}
N/A

\section{Configuring SSFM}
Once the complete package has been added to your project, an additional module RTA-SUM SSFM should appear.
From here you can open and start to configure the module.
Specific descriptions of each configurable parameter available are available in the “Description” view of 
\ifnum \DocumentGroup=1
{
RTA-CAR ISOLAR-B
}
\else
{
ISOLAR-A with BCT add-on 
}\fi when a parameter is selected.

\section{Generating Artifacts}
Generating the .c/.h/.arxml files from your configuration is similar to other BSW modules delivered from ETAS.
The Code Generation perspective of the 
\ifnum \DocumentGroup=1
{
RTA-CAR ISOLAR-B
}
\else
{
ISOLAR-A with BCT add-on 
}\fi , should present the user with a dialog to select specific module within the project to generate.
\ifnum \DocumentGroup=1
{
%Nothing
}
\else
{
Either a Partial Build or Full Build can be used, depending on the needs of the user.
}\fi
In case of any configuration errors present in SSFM, the console window will output a description of any errors found during generation.
Further, when the user is configuring the module, the Problems Log will show any active errors in the configuration of SSFM.
A third visual aid is the appearance of red X’s in the configuration view over items which are invalid.
Simply hovering over these X’s will provide the user a description of the problem.
These are static validations.

\section{Configuring the RTE}
After the user has configured and generated the artifacts for SSFM, the output SWCD file \linebreak(\file{SSFM\_SWCD.arxml}) can be used for the RTE configuration.
This is a complete SWC description per the AUTOSAR standard of all the users configured ports and interfaces required for the project specific configuration.
This SWCD is a standard ASW component and not a module description.
Therefore, there are no special considerations to be taken into account when configuring the RTE.
It is recommended that if specific changes are required to the SWCD, the user not manually edit this file, rather update and regenerate the configuration with the required changes.

Note: RTE command used for RTE generation

\begin{rtecommandline}
--strict-unconnected-rport-check=off --samples=memmap -eao=disable --os-define-osenv=RTAOS40 -dv8aur=1
\end{rtecommandline}